\documentclass[aspectratio=169, xcolor=table]{beamer}
% PACKAGES & CONFIGURATION
\usepackage[english]{babel}
\usepackage[utf8]{inputenc}
\usepackage{graphicx}
\usepackage{booktabs}
\usepackage{hyperref}
\usepackage{xcolor}
\usepackage{listings}
\usepackage{lstautogobble}
\usepackage{fancyvrb}

% THEME & COLORS
\mode<presentation>
{
\usetheme{default}
\usecolortheme{default}

% Define custom tech-forward color palette
\definecolor{PrimaryDark}{HTML}{1e293b}      % Slate 900 - Deep dark for backgrounds
\definecolor{PrimaryAccent}{HTML}{0ea5e9}    % Sky Blue - Modern accent
\definecolor{SecondaryAccent}{HTML}{10b981}  % Emerald - Secondary accent
\definecolor{TextLight}{HTML}{f1f5f9}        % Slate 100 - Light text
\definecolor{TextDark}{HTML}{1e293b}         % Slate 900 - Dark text
\definecolor{CodeBg}{HTML}{0f172a}           % Slate 950 - Code background
\definecolor{CodeText}{HTML}{e2e8f0}         % Slate 200 - Code text
\definecolor{AlertColor}{HTML}{f97316}       % Orange - For emphasis

% Set overall colors
\setbeamercolor{background canvas}{bg=PrimaryDark}
\setbeamercolor{normal text}{fg=TextLight}
\setbeamercolor{title}{fg=PrimaryAccent, bg=}
\setbeamercolor{subtitle}{fg=SecondaryAccent}
\setbeamercolor{structure}{fg=PrimaryAccent}
\setbeamercolor{alerted text}{fg=AlertColor}
\setbeamercolor{itemize item}{fg=PrimaryAccent}
\setbeamercolor{itemize subitem}{fg=SecondaryAccent}
\setbeamercolor{enumerate item}{fg=PrimaryAccent}
\setbeamercolor{enumerate subitem}{fg=SecondaryAccent}
\setbeamercolor{block title}{fg=TextLight, bg=PrimaryAccent}
\setbeamercolor{block body}{fg=TextLight, bg=}
\setbeamercolor{footline}{fg=TextLight, bg=}
\setbeamercolor{frametitle}{fg=PrimaryAccent}

% Set fonts
\usefonttheme{structurebold}
\setbeamerfont{frametitle}{size=\large, series=\bfseries}
\setbeamerfont{title}{size=\Huge, series=\bfseries}
\setbeamerfont{subtitle}{size=\Large}
\setbeamerfont{normal text}{size=\normalsize}
\setbeamerfont{itemize/enumerate body}{size=\large}
\setbeamerfont{itemize/enumerate subbody}{size=\normalsize}

% Customize templates
\setbeamertemplate{navigation symbols}{}  % Remove navigation symbols
\setbeamertemplate{caption}[numbered]

% Custom footer with slide number and author
\setbeamertemplate{footline}{
	\leavevmode%
	\hbox{%
		\begin{beamercolorbox}[wd=.5\paperwidth,ht=2.25ex,dp=1ex,left]{author in head/foot}%
			\hspace{1ex}\insertshortauthor
		\end{beamercolorbox}%
		\begin{beamercolorbox}[wd=.5\paperwidth,ht=2.25ex,dp=1ex,right]{date in head/foot}%
			\insertframenumber{}\hspace{1ex}
		\end{beamercolorbox}}%
	\vskip0pt%
}

% Customize frame title style
\setbeamertemplate{frametitle}{
\vskip0.5em
{\usebeamerfont{frametitle}\usebeamercolor[fg]{frametitle}\insertframetitle}
\vskip-0.5em
{\color{gray}\rule{\linewidth}{0.4pt}}
}
}

% CODE HIGHLIGHTING

% Terraform & HashiCorp Configuration Language (HCL)
\lstdefinelanguage{Terraform}{
	sensitive=true,
	morekeywords={
		terraform,required_providers,required_version,provider,resource,data,module,
		variable,locals,output,backend,cloud,workspace,lifecycle,depends_on,
		count,for_each,dynamic,provisioner,connection,validation,precondition,
		postcondition,check,assert,source,version,type,description,default,nullable,
		tags,bucket,region,rule,status,id,merge
	},
	morekeywords=[2]{true,false,null},
	morecomment=[l]{\#},
	morecomment=[l]{//},
	morecomment=[s]{/*}{*/},
	morestring=[b]"
}

\lstset{
	autogobble=true,
	basicstyle=\ttfamily\small\color{CodeText},
	backgroundcolor=\color{CodeBg},
	breaklines=true,
	captionpos=b,
	commentstyle=\color{gray},
	keywordstyle=\color{PrimaryAccent}\bfseries,
	numberstyle=\tiny\color{gray},
	stringstyle=\color{SecondaryAccent},
	showstringspaces=false,
	frame=single,
	frameround=tttt,
	framerule=0pt,
	rulecolor=\color{gray},
	xleftmargin=1em,
	xrightmargin=1em,
	framexleftmargin=0.8em,
	framexrightmargin=0.8em,
}


\lstset{
  language=Terraform,
  keywordstyle=[2]\color{SecondaryAccent}\bfseries
}

\title{Terraform, from First Principles}
\subtitle{The practical fundamentals of infrastructure as code}
\author{Marcel Barlik}
\date{18 August 2026}

\begin{document}

\begin{frame}[plain]
  \titlepage
\end{frame}

\begin{frame}{Today: 30 Minutes to a Safe First Workflow}
  \begin{itemize}
    \item \textbf{What Terraform is} (5 min) --- infrastructure as reviewed, repeatable code
    \item \textbf{Define the contracts} (8 min) --- providers, resources, inputs and locals
    \item \textbf{Use the workflow} (8 min) --- init, fmt, validate, plan and apply
    \item \textbf{Work safely as a team} (5 min) --- state and review
  \end{itemize}

  \vspace{1em}
  \begin{alertblock}{Outcome}
    You will be able to read a Terraform configuration, make a small infrastructure change,
    review its plan, and know when \emph{not} to apply it.
  \end{alertblock}
\end{frame}

\section{Why Terraform?}

\begin{frame}{What Is Terraform?}
  \begin{block}{Definition}
    Terraform is an \textbf{Infrastructure as Code (IaC)} tool. The infrastructure and
    configuration defined in files is the source of truth. Terraform compares that desired
    state with the real world, and based on this it proposes the necessary changes.
  \end{block}

  \vspace{0.7em}
  \begin{columns}[T]
    \column{0.48\textwidth}
    \textbf{It manages}
    \begin{itemize}
      \item Cloud resources: networks, compute, storage, databases
      \item SaaS configuration: dashboards, teams, access controls
      \item Internal platforms exposed through APIs
    \end{itemize}

    \column{0.48\textwidth}
    \textbf{It gives you}
    \begin{itemize}
      \item A version-controlled description of intent
      \item A preview of changes before they happen
      \item Repeatable environments and a change history
    \end{itemize}
  \end{columns}
\end{frame}

\begin{frame}{Manual Configuration vs.\ Terraform}
  \begin{columns}[T]
    \column{0.49\textwidth}
    \textbf{Clicking in a UI}
    \begin{itemize}
      \item Easy to make an unrecorded change
      \item Hard to reproduce in another environment
      \item Review happens after the fact
      \item Restoring a previous layout is manual
    \end{itemize}

    \column{0.49\textwidth}
    \textbf{Terraform}
    \begin{itemize}
      \item Desired infrastructure lives in Git
      \item One configuration can target each environment
      \item \texttt{plan} shows the proposed change first
      \item History, review and rollback use normal engineering practice
    \end{itemize}
  \end{columns}

  \vspace{1em}
  \textcolor{SecondaryAccent}{\textbf{Terraform does not replace cloud, SaaS, or internal platform.}}
  It makes the desired configuration explicit and repeatable.
\end{frame}

\begin{frame}{Terraform vs.\ a Script}
  \begin{columns}[T]
    \column{0.48\textwidth}
    \textbf{Imperative: describe the steps}
    \begin{block}{Intent in a script}
      ``Create a bucket; add these tags; then create a policy.''
    \end{block}
    \begin{itemize}
      \item The author manages ordering and retries
      \item Re-running safely is extra work
      \item The script does not inherently know what already exists
    \end{itemize}

    \column{0.48\textwidth}
    \textbf{Declarative: describe the outcome}
    \begin{block}{Intent in Terraform}
      ``This bucket, its tags, and its policy should exist.''
    \end{block}
    \begin{itemize}
      \item Terraform builds the dependency graph
      \item \texttt{plan} calculates the delta from current state
      \item \texttt{apply} changes only what is needed
    \end{itemize}
  \end{columns}
\end{frame}

\begin{frame}{Terraform's Change Flow}
  \begin{center}
    \Large \textcolor{PrimaryAccent}{Configure} $\longrightarrow$ \textcolor{SecondaryAccent}{Plan} $\longrightarrow$ \textcolor{AlertColor}{Apply}
  \end{center}

  \vspace{1em}
  \begin{columns}[T]
    \column{0.31\textwidth}
    \textbf{Configure}
    \begin{itemize}
      \item Describe the desired outcome
      \item Store it in version control
    \end{itemize}

    \column{0.31\textwidth}
    \textbf{Plan}
    \begin{itemize}
      \item Preview the proposed changes
      \item Review before anything changes
    \end{itemize}

    \column{0.31\textwidth}
    \textbf{Apply}
    \begin{itemize}
      \item Make the approved changes
      \item Verify the resulting infrastructure
    \end{itemize}
  \end{columns}

  \vspace{0.8em}
  \begin{alertblock}{Important boundary}
    \texttt{plan} is read-only. \texttt{apply} is the explicit step that changes remote infrastructure.
  \end{alertblock}
\end{frame}

\begin{frame}{Terraform's Mental Model}
  Terraform calculates the plan by bringing together desired configuration, its memory of
  managed objects, and the platform it is managing.

  \vspace{0.7em}
  \begin{columns}[T]
    \column{0.32\textwidth}
    \textbf{Configuration}
    \begin{itemize}
      \item \texttt{.tf} files
      \item Desired infrastructure
      \item Inputs and rules
    \end{itemize}

    \column{0.32\textwidth}
    \textbf{State}
    \begin{itemize}
      \item Terraform's record
      \item Maps code to remote objects
      \item Enables diffing
    \end{itemize}

    \column{0.32\textwidth}
    \textbf{Provider APIs}
    \begin{itemize}
      \item Cloud, SaaS, or internal APIs
      \item Reads actual state
      \item Creates or updates objects
    \end{itemize}
  \end{columns}

  \vspace{0.8em}
  \begin{block}{The key idea}
    Terraform converges the remote system toward the configuration. It is not a script
    that blindly re-runs every line.
  \end{block}
\end{frame}

\section{Define the contracts}

\begin{frame}{Start with Four Building Blocks}
  Most Terraform configuration is easier to read once four words are clear.

  \vspace{0.8em}
  \begin{columns}[T]
    \column{0.48\textwidth}
    \begin{block}{Provider}
      The plugin that speaks to a platform's API.
    \end{block}
    \begin{block}{Resource}
      A remote object Terraform is responsible for.
    \end{block}

    \column{0.48\textwidth}
    \begin{block}{Variable}
      A value supplied from outside the module.
    \end{block}
    \begin{block}{Local value}
      A value calculated and named inside the module.
    \end{block}
  \end{columns}

  \vspace{2em}
  We will build one small configuration using each block in that order.
\end{frame}

\begin{frame}{A Provider Is Terraform's API Adapter}
  Terraform Core understands configuration, state and the plan. A \textbf{provider}
  understands a particular remote system.

  \vspace{0.8em}
  \begin{center}
    \texttt{Terraform Core}
    \quad $\longleftrightarrow$ \quad
    \textcolor{SecondaryAccent}{\texttt{AWS provider}}
    \quad $\longleftrightarrow$ \quad
    \texttt{AWS APIs}
  \end{center}

  \vspace{1em}
  The provider defines:
  \begin{itemize}
    \item which \textbf{resource} types Terraform can manage;
    \item which settings those \textbf{resources} accept; and
    \item how to read, create, update and delete them through the API.
  \end{itemize}

  \begin{block}{Examples}
    AWS, Azure, GitHub, Datadog and many internal platforms can all have providers.
  \end{block}
\end{frame}

\begin{frame}[fragile]{A Provider Has Two Separate Contracts}
  \begin{columns}[T]
    \column{0.49\textwidth}
    \textbf{1. Require the plugin}
    \begin{lstlisting}[basicstyle=\ttfamily\tiny\color{CodeText}]
terraform {
  required_providers {
    aws = {
      source  = "hashicorp/aws"
      version = "~> 5.0"
    }
  }
}
    \end{lstlisting}
    \small Which provider may Terraform install?

    \column{0.49\textwidth}
    \textbf{2. Configure the connection}
    \begin{lstlisting}[basicstyle=\ttfamily\tiny\color{CodeText}]
provider "aws" {
  region = var.aws_region
}
    \end{lstlisting}
    \small How should this module use that provider?
  \end{columns}

  \vspace{0.8em}
  \begin{alertblock}{Keep identity outside the code}
    Region and other non-secret settings may be configuration. Credentials should come
    from the platform's normal environment, profile or workload identity mechanism.
  \end{alertblock}
\end{frame}

\begin{frame}{What Happens When You Run \texttt{terraform init}?}
  \texttt{init} prepares the current module; it does not create infrastructure.

  \vspace{0.8em}
  \begin{enumerate}
    \item Terraform reads the \texttt{required\_providers} block.
    \item It downloads a compatible provider plugin.
    \item It records the exact selected version in \texttt{.terraform.lock.hcl}.
  \end{enumerate}

  \vspace{0.8em}
  \begin{block}{Team habit}
    Commit the lock file. Everyone and CI then start from the same provider selection.
  \end{block}

  \vspace{0.8em}
  \textcolor{SecondaryAccent}{\textbf{So far:}} Terraform can talk to AWS, but we have not
  asked it to manage anything.
\end{frame}

\begin{frame}[fragile]{A Resource Is an Object Terraform Manages}
  \begin{lstlisting}[basicstyle=\ttfamily\scriptsize\color{CodeText}]
resource "aws_s3_bucket" "assets" {
  bucket = "payments-staging-assets"
}
  \end{lstlisting}

  \vspace{0.5em}
  \begin{itemize}
    \item \textcolor{PrimaryAccent}{\texttt{resource}} declares an object Terraform will manage.
    \item \textcolor{SecondaryAccent}{\texttt{aws\_s3\_bucket}} is the resource type, defined by the provider.
    \item \textcolor{SecondaryAccent}{\texttt{assets}} is the local name, chosen by us.
    \item Together, the type and name form the resource's address:
    \textcolor{SecondaryAccent}{\texttt{aws\_s3\_bucket.assets}}.
  \end{itemize}
\end{frame}

\begin{frame}[fragile]{Inside a Block: Arguments and Nested Blocks}
  \begin{columns}[T]
    \column{0.48\textwidth}
    \textbf{Argument}
    \begin{lstlisting}[basicstyle=\ttfamily\scriptsize\color{CodeText}]
bucket = "payments-assets"
    \end{lstlisting}
    Assigns an expression to a named setting.

    \column{0.48\textwidth}
    \textbf{Nested block}
    \begin{lstlisting}[basicstyle=\ttfamily\scriptsize\color{CodeText}]
rule {
  status = "Enabled"
}
    \end{lstlisting}
    Groups related settings into a structure.
  \end{columns}

  \vspace{1em}
  \begin{block}{The provider documentation is the contract}
    It tells us which arguments are required, which blocks are valid, and which attributes
    become available after creation.
  \end{block}
\end{frame}

\begin{frame}[fragile]{References Connect Managed Objects}
  \begin{lstlisting}[basicstyle=\ttfamily\scriptsize\color{CodeText}]
resource "aws_s3_bucket_lifecycle_configuration" "assets" {
  bucket = aws_s3_bucket.assets.id

  rule {
    id     = "archive-old-assets"
    status = "Enabled"
  }
}
  \end{lstlisting}

  \textcolor{SecondaryAccent}{\texttt{aws\_s3\_bucket.assets.id}} means:
  \begin{itemize}
    \item find the resource at \textcolor{SecondaryAccent}{\texttt{aws\_s3\_bucket.assets}};
    \item read its resulting \textcolor{SecondaryAccent}{\texttt{id}};
    \item use that value here.
  \end{itemize}

  \vspace{0.8em}
  \textcolor{SecondaryAccent}{\textbf{Terraform also learns the order:}}
  create or read the bucket before configuring its lifecycle.
\end{frame}

\begin{frame}[fragile]{A Variable Is an Input to the Module}
  The declaration defines what the module accepts:
  \begin{lstlisting}[basicstyle=\ttfamily\scriptsize\color{CodeText}]
variable "environment" {
  description = "Environment this bucket belongs to"
  type        = string
  nullable    = false
}
  \end{lstlisting}

  The resource reads that input as \texttt{var.environment}:
  \begin{lstlisting}[basicstyle=\ttfamily\scriptsize\color{CodeText}]
bucket = "payments-${var.environment}-assets"
  \end{lstlisting}

  \begin{block}{Why declare a type?}
    Terraform can reject an invalid input before it attempts a remote change.
  \end{block}
\end{frame}

\begin{frame}[fragile]{Declaring a Variable Is Not Assigning Its Value}
  \begin{columns}[T]
    \column{0.48\textwidth}
    \textbf{The module declares}
    \begin{lstlisting}[basicstyle=\ttfamily\scriptsize\color{CodeText}]
variable "environment" {
  type = string
}
    \end{lstlisting}
    Usually stored in \textcolor{SecondaryAccent}{\texttt{variables.tf}}.

    \column{0.48\textwidth}
    \textbf{The caller supplies}
    \begin{lstlisting}[basicstyle=\ttfamily\scriptsize\color{CodeText}]
environment = "staging"
    \end{lstlisting}
    For example, in \textcolor{SecondaryAccent}{\texttt{staging.tfvars}}.
  \end{columns}

  \vspace{1em}
  A team can instead inject values through CI or \texttt{TF\_VAR\_...}.

  \begin{alertblock}{Security note}
  Never put secrets in a committed \textcolor{SecondaryAccent}{\texttt{.tfvars}} file.
  \end{alertblock}
\end{frame}

\begin{frame}[fragile]{A Local Names a Value Calculated Here}
  \begin{lstlisting}[basicstyle=\ttfamily\scriptsize\color{CodeText}]
locals {
  resource_name = "payments-${var.environment}"
  common_tags = {
    Environment = var.environment
    ManagedBy   = "terraform"
  }
}
  \end{lstlisting}

  \begin{columns}[T]
    \column{0.48\textwidth}
    \textbf{Variable: \textcolor{SecondaryAccent}{\texttt{var.environment}}}

    Supplied by the caller; part of the module's interface.

    \column{0.48\textwidth}
    \textbf{Local: \textcolor{SecondaryAccent}{\texttt{local.resource\_name}}}

    Derived inside the module; removes repetition and gives an expression a useful name.
  \end{columns}
\end{frame}

\begin{frame}[fragile]{Now Read the Whole Value Flow}
  \begin{lstlisting}[basicstyle=\ttfamily\scriptsize\color{CodeText}]
locals {
  resource_name = "payments-${var.environment}"
}
resource "aws_s3_bucket" "assets" {
  bucket = "${local.resource_name}-assets"
}
  \end{lstlisting}

  \vspace{0.5em}
  \begin{columns}[T]
    \column{0.31\textwidth}
    \textcolor{PrimaryAccent}{\textbf{1. Input variable}}

    \textcolor{SecondaryAccent}{\texttt{var.environment}}

    \small Supplied as \textcolor{SecondaryAccent}{\texttt{staging}}.

    \column{0.31\textwidth}
    \textcolor{PrimaryAccent}{\textbf{2. Calculated local}}

    \textcolor{SecondaryAccent}{\texttt{local.resource\_name}}

    \small Becomes \textcolor{SecondaryAccent}{\texttt{payments-staging}}.

    \column{0.31\textwidth}
    \textcolor{PrimaryAccent}{\textbf{3. Resource argument}}

    \texttt{bucket}

    \small Receives
    \textcolor{SecondaryAccent}{\texttt{payments-staging-assets}}.
  \end{columns}

  \vspace{1.6em}
  \textcolor{SecondaryAccent}{\textbf{The idea:}} one input can feed a local, and that local
  can then feed many resource arguments.
\end{frame}

\begin{frame}{Files Are for Humans; the Directory Is the Module}
  Terraform loads all \texttt{.tf} files in one directory together as a \textbf{module}.

  \vspace{0.6em}
  \begin{center}
    \small
    \renewcommand{\arraystretch}{1.3}
    \arrayrulecolor{gray}
    \begin{tabular}{|p{0.31\textwidth}|p{0.55\textwidth}|}
      \hline
      \rowcolor{PrimaryAccent}
      \textbf{File} & \textbf{Typical purpose} \\
      \hline
      \texttt{versions.tf} & Provider requirements \\
      \hline
      \texttt{providers.tf} & Provider configuration \\
      \hline
      \texttt{variables.tf} & Input contracts \\
      \hline
      \texttt{main.tf} & Resources and local values \\
      \hline
      \texttt{outputs.tf} & Values exposed to callers \\
      \hline
    \end{tabular}
  \end{center}

  \vspace{0.5em}
  \textcolor{SecondaryAccent}{\textbf{Important:}} file order does not control creation
  order. References between values and resources do.
\end{frame}

\section{Use the workflow}

\begin{frame}[fragile]{Terraform Commands: The Two Halves}
  \begin{lstlisting}[basicstyle=\ttfamily\scriptsize\color{CodeText}]
terraform init

terraform fmt
terraform validate

terraform plan
terraform apply
  \end{lstlisting}

  \begin{columns}[T]
    \column{0.48\textwidth}
    \textbf{Prepare and check}

    Work on the local configuration.

    \column{0.48\textwidth}
    \textbf{Preview and change}

    Interact with the target platform.
  \end{columns}

  \vspace{2em}
  Run these commands from the module directory. \texttt{init} is also needed when provider
  requirements or the state backend change.
\end{frame}

\begin{frame}{\texttt{fmt} and \texttt{validate} Answer Different Questions}
  \begin{columns}[T]
    \column{0.48\textwidth}
    \begin{block}{\texttt{terraform fmt}}
      ``Is the configuration formatted consistently?''
    \end{block}

    It rewrites files into Terraform's standard style.

    \column{0.48\textwidth}
    \begin{block}{\texttt{terraform validate}}
      ``Is the configuration internally valid?''
    \end{block}

    It checks syntax, types and references without proposing a remote change.
  \end{columns}

  \vspace{1em}
  \textcolor{SecondaryAccent}{\textbf{Both are fast feedback.}} Neither proves that the
  change is safe for a particular account or environment; that is the plan's job.
\end{frame}

\begin{frame}{A Plan Is Terraform's Proposed Change}
  \texttt{terraform plan} brings together three views:

  \vspace{0.8em}
  \begin{center}
    desired configuration
    \quad + \quad
    Terraform state
    \quad + \quad
    current provider API data
  \end{center}

  \vspace{1em}
  The result is a preview of what \emph{would} change. Planning may read remote systems
  and require credentials, but it does not apply the proposed infrastructure changes.

  \begin{alertblock}{A successful plan is not approval}
    ``Terraform can do this'' is different from ``we intend to do this.'' Read the result.
  \end{alertblock}
\end{frame}

\begin{frame}{Read the Plan as a Code Review}
  Focus first on Terraform's action symbols:

  \vspace{0.5em}
  \begin{center}
    \begin{tabular}{c l p{0.50\textwidth}}
      \toprule
      \textbf{Symbol} & \textbf{Action} & \textbf{Question to ask} \\
      \midrule
      \textcolor{SecondaryAccent}{\texttt{+}} & create & Is this object expected and named correctly? \\
      \textcolor{AlertColor}{\texttt{-}} & destroy & Why is it being removed? Is this expected?  \\
      \textcolor{PrimaryAccent}{\texttt{\textasciitilde}} & update & Is each changed field intentional? \\
      \textcolor{AlertColor}{\texttt{-/+}} & replace & What is being replaced? Is potential downtime/data loss acceptable? \\
      \bottomrule
    \end{tabular}
  \end{center}

  \vspace{0.5em}
  \begin{alertblock}{Stop condition}
    Do not apply an unexplained destroy or replacement. The plan is the last inexpensive
    place to catch a wrong value, wrong target account, or typo.
  \end{alertblock}
\end{frame}

\begin{frame}{Apply Only the Plan You Understand}
  A safe review answers four questions:

  \vspace{0.5em}
  \begin{enumerate}
    \item Am I targeting the intended account and environment?
    \item Are the resource addresses the ones I expected?
    \item Can I explain every create, update, replacement and destroy?
    \item Is the operational impact acceptable?
  \end{enumerate}

  \vspace{0.8em}
  \begin{block}{The boundary}
    \texttt{apply} is the explicit command that asks the provider to change remote objects.
  \end{block}
\end{frame}

\begin{frame}{State Is Terraform's Memory}
  \textbf{State} records which remote object belongs to each resource address.

  \vspace{0.4em}
  \begin{columns}[T]
    \column{0.48\textwidth}
    \textbf{Configuration says}
    \begin{itemize}
      \item ``Manage a bucket at this address.''
      \item ``These settings are desired.''
    \end{itemize}

    \column{0.48\textwidth}
    \textbf{State remembers}
    \begin{itemize}
      \item ``That address maps to bucket ID 123.''
      \item Provider metadata needed to manage it.
    \end{itemize}
  \end{columns}

  \vspace{0.8em}
  \begin{alertblock}{Do not delete state to ``start over''}
    Removing Terraform's mapping does not remove the remote object; it makes Terraform
    forget that the object is already managed.
  \end{alertblock}
\end{frame}

\begin{frame}{Drift Is a Difference Terraform Discovers}
  \begin{center}
    \small
    configuration says A
    \quad $\neq$ \quad
    provider API reports B
  \end{center}

  \vspace{0.8em}
  \textbf{Drift} often appears when someone changes a managed object directly in its UI or API.

  \vspace{0.6em}
  \begin{enumerate}
    \item A plan reads the remote object and reveals the difference.
    \item The team decides which version is intended.
    \item Code and infrastructure are reconciled through a reviewed change.
  \end{enumerate}

  \vspace{0.6em}
  \textcolor{SecondaryAccent}{\textbf{Habit:}} use Terraform as the normal route for changing
  objects Terraform manages.
\end{frame}

\section{Work safely as a team}

\begin{frame}{Shared Work Needs Shared State}
  Local \texttt{terraform.tfstate} is suitable for learning, not for a team managing the
  same infrastructure.

  \vspace{0.6em}
  \begin{block}{Remote backend}
    Stores state in a shared, durable location so engineers and CI use the same memory.
  \end{block}

  \begin{block}{Locking}
    Prevents two writers from changing the same state at the same time.
  \end{block}

  \begin{alertblock}{Treat state as sensitive operational data}
    It can contain resource details and sometimes secret values. Encrypt it, restrict access,
    keep backups, and never commit it to Git.
  \end{alertblock}
\end{frame}

\begin{frame}{A Safe Team Change Has One Clear Path}
  \begin{center}
    branch
    $\longrightarrow$ check
    $\longrightarrow$ plan
    $\longrightarrow$ review
    $\longrightarrow$ apply
    $\longrightarrow$ verify
  \end{center}

  \vspace{0.8em}
  \begin{itemize}
    \item Keep the code change focused; separate resource changes from provider upgrades.
    \item Review the plan for the intended environment in the pull request.
    \item Let CI or a designated operator apply with environment-scoped credentials.
    \item Verify the remote result after apply.
  \end{itemize}

  \begin{block}{The invariant}
    The code reviewed, the values planned and the infrastructure applied must describe the
    same change.
  \end{block}
\end{frame}

\begin{frame}{Three More Terms You Will Meet}
  \begin{description}
    \item[Output] A value a module exposes to its caller, such as a bucket ID.
    \item[Data source] A read-only lookup of an object Terraform does not own here.
    \item[Child module] A reusable Terraform package called by another module.
  \end{description}

  \vspace{0.8em}
  \begin{block}{A sensible order for learning}
    Start with resources, variables and local values in one module. Add outputs when callers
    need results, data sources when ownership is elsewhere, and child modules when reuse is
    useful -- just like when you program and feel the need to refactor.
  \end{block}
\end{frame}

\begin{frame}{Pause When the Plan Surprises You}
  Do not apply yet if you see:

  \vspace{0.5em}
  \begin{itemize}
    \item an unexplained destroy or replacement;
    \item resources from the wrong account or environment;
    \item unrelated changes you did not intend; or
    \item a provider or lock-file upgrade mixed into an ordinary resource change.
  \end{itemize}

  \vspace{0.7em}
  Check the selected values, credentials, provider configuration, resource references and
  state before trying again.

  \begin{alertblock}{Recovery tools are not everyday shortcuts}
    Commands such as \texttt{-target}, import and direct state manipulation require a specific,
    understood recovery plan.
  \end{alertblock}
\end{frame}

\begin{frame}{What You Should Now Be Able to Do}
  \begin{enumerate}
    \item Explain the job of a provider and the difference between requiring and configuring it.
    \item Read a resource, its arguments, nested blocks, variables, locals and references.
    \item Follow \texttt{init} $\rightarrow$ \texttt{fmt} $\rightarrow$ \texttt{validate}
      $\rightarrow$ \texttt{plan} $\rightarrow$ \texttt{apply}.
    \item Treat the plan as a change review and state as sensitive shared operational data.
    \item Recognise when a proposed change is safe to apply---and when to stop.
  \end{enumerate}

  \vspace{1em}
  \begin{alertblock}{One principle to remember}
    \textbf{Never apply a Terraform plan you cannot explain.}
  \end{alertblock}
\end{frame}

\begin{frame}{Questions}
  \begin{center}
    \vspace{1.5em}
    \Huge \textcolor{PrimaryAccent}{Questions?}
    \vspace{1.5em}

    {\large Suggested next step: make one safe, non-production Terraform change.}
    \vspace{0.8em}

    {\small \textcolor{gray}{Start with any small module; plan before you apply.}}
  \end{center}
\end{frame}

\appendix

\begin{frame}[allowframebreaks]{References}
  \small
  \begin{itemize}
    \item Terraform documentation: \url{https://developer.hashicorp.com/terraform/docs}
    \item Terraform CLI reference: \url{https://developer.hashicorp.com/terraform/cli}
    \item Terraform Registry: \url{https://registry.terraform.io/}
  \end{itemize}
\end{frame}

\end{document}
